problem:  
Let \( N_{k,l} = \mathrm{SU}(3)/\mathrm{U}(1)_{k,l} \) be an Aloff–Wallach space with the standard isotropy decomposition \( \mathfrak{m} = \mathfrak{m}_1 \oplus \mathfrak{m}_2 \oplus \mathfrak{m}_3 \) into three irreducible \(\mathrm{Ad}(\mathrm{U}(1)_{k,l})\)-modules. Every \(\mathrm{SU}(3)\)-invariant metric can be written as  
\[
g_{\lambda_1, \lambda_2, \lambda_3} = \lambda_1^2 Q|_{\mathfrak{m}_1} + \lambda_2^2 Q|_{\mathfrak{m}_2} + \lambda_3^2 Q|_{\mathfrak{m}_3},
\]
where \(Q\) is the negative Killing form on \(\mathfrak{su}(3)\) and each \(\lambda_i > 0\).  

The natural homogeneous fibration  
\[
S^3 \longrightarrow N_{k,l} \xrightarrow{\pi} \mathbb{C}P^2
\]
suggests that as one of the parameters \(\lambda_i\) tends to zero, the metric may collapse (after appropriate rescaling) to a metric on the base or the fiber with possible Einstein properties.  

**Problem:**  
Investigate whether there exist sequences of \(\mathrm{SU}(3)\)-invariant metrics \(g_{\lambda_1(t), \lambda_2(t), \lambda_3(t)}\) on \(N_{k,l}\) such that  

1. The metrics have *nonnegative sectional curvature* for all \(t > 0\);  
2. One of the scale factors, say \(\lambda_3(t)\), tends to \(0\) as \(t \to 0\);  
3. The rescaled metrics  
   \[
   \tilde g_t = \lambda_3(t)^{-2} g_{\lambda_1(t), \lambda_2(t), \lambda_3(t)}
   \]
   converge in the Cheeger–Gromov sense to a smooth homogeneous Einstein metric on the submanifold corresponding to the collapsed fiber (e.g. an \(\mathrm{SU}(2)\) orbit or a Berger-type 3-sphere).

Equivalently, does there exist a nonnegatively curved invariant metric deformation on an Aloff–Wallach space that collapses smoothly—via the homogeneous fibration structure—to a lower-dimensional *Einstein* homogeneous space?  

If such a path exists, describe the limiting curvature distribution and determine whether the total scalar curvature (appropriately normalized by volume) remains bounded along the deformation.

Why is it a "good" problem:  
This formulation gives a concrete and logically sound geometric degeneration problem: it asks whether **nonnegative curvature can be preserved under a controlled homogeneous fibration collapse** on Aloff–Wallach spaces, with an Einstein limit appearing on the fiber or base after suitable rescaling. The setting is deeply natural because Aloff–Wallach spaces arise precisely as homogeneous fibrations \(S^3 \to N_{k,l} \to \mathbb{C}P^2\), and the question probes the subtle relationship between **nonnegative curvature**, **collapse theory**, and **Einstein geometry** within this class.  
Resolving the problem would require understanding the fine curvature inequalities for invariant metrics, the analytical behavior of the sectional curvature under the fibration collapse, and the rescaling mechanisms that preserve homogeneity in limits. This bridges **homogeneous Riemannian geometry**, **Cheeger–Gromov collapse analysis**, and **Einstein metric theory**, embodying the hallmark of a "good" problem: it is simple to state, rooted in canonical structures, yet its resolution demands a synthesis of deep geometric and analytic insights.